\chapter*{Kurzfassung}

Diese Bachelorarbeit beschreibt den Entwurf und die prototypische 
Entwicklung des Frontends der Webanwendung \textit{Bazario}, 
einer Multi-Rollen-Plattform, die E-Commerce und 
Dienstleistungsbuchungen in einem einheitlichen System vereint.

Die Plattform unterscheidet vier Nutzerrollen: Käufer, Verkäufer,
Dienstleister und Administratoren. Das entwickelte Frontend deckt
die Funktionsbereiche der ersten drei Rollen ab; die
administrativen Funktionen erfolgen über das vom Backend
bereitgestellte \textit{Admin}-Dashboard. React bildet das Kernframework, React Router übernimmt das
clientseitige Routing mit rollenbasiertem Seitenschutz und
die Bibliothek Zustand verwaltet den globalen Anwendungszustand
für Authentifizierung, Nutzerrolle und Warenkorb. Die Zahlungsabwicklung erfolgt
über Stripe, die Authentifizierung erfolgt tokenbasiert über
Laravel Sanctum mit Bearer Tokens.

Neben klassischen E-Commerce-Funktionen wie Produktkatalog, 
Warenkorb und Checkout umfasst \textit{Bazario} ein 
Buchungssystem für Dienstleistungen, ein integriertes 
Messaging-System sowie ein Werbe- und Ankündigungsmodul. 
Die Qualität der Umsetzung wird durch Unit-, Integrations- 
und End-to-End-Tests abgesichert und anhand von 
Usability-Kriterien bewertet.

\textbf{Schlüsselwörter:} E-Commerce, Marktplatz,
Multi-Rollen-System, React, Single-Page Application,
Zustand, Stripe, Laravel Sanctum, Rollenbasierte
Zugriffskontrolle, Frontend-Entwicklung